% !TEX root = 00_instructivo.tex
\chapter{Transductores electroquímicos y MEMS: Calidad de agua y aire}

\section*{Objetivos}
\begin{itemize}
    \item Medir sólidos disueltos totales (TDS) en agua usando un sensor de conductividad.
    \item Determinar el pH de muestras de agua (ácido, neutro, básico).
    \item Detectar gases contaminantes (CO, metano, etanol) en el aire usando un sensor MEMS MiCS-5524.
    \item Integrar los tres sensores para generar un perfil básico de contaminación ambiental.
\end{itemize}

\section*{Preguntas previas}
\begin{enumerate}
    \item ¿Qué es el TDS y cómo se relaciona con la conductividad eléctrica del agua?
    \item ¿Por qué es importante compensar la temperatura en las mediciones de pH y TDS?

> Mía:
\item ¿Qué tipo de gases puede detectar el MiCS-5524? ¿Cuál es el principio de funcionamiento de un sensor de óxido metálico (MOS)?
    \item Investigue: ¿Cuáles son los límites máximos permisibles de TDS y pH para agua potable según la OMS?
\end{enumerate}

\section*{Materiales y equipo}
\begin{itemize}
    \item Arduino Uno.
    \item Sensor de TDS (Gravity: Analog TDS Sensor).
    \item Sensor de pH (Gravity: Analog pH Meter Kit).
    \item Sensor de gas MiCS-5524 (Fermion MEMS).
    \item Muestras de agua: destilada, potable, agua con sal, agua con vinagre (ácido), agua con bicarbonato (básico).
    \item Soluciones buffer de calibración (pH 4.0, pH 7.0) y solución de calibración TDS (1413 µS/cm).
    \item Recipientes plásticos pequeños (vasos de precipitado).
    \item Mini protoboard y cables Dupont.
    \item Fuente de aire contaminado (encendedor sin llama para gas, alcohol en algodón).
\end{itemize}

\section*{Procedimiento}
\begin{enumerate}
    \item \textbf{Calibración de sensores (fundamental):}
    \begin{enumerate}
        \item \textbf{pH:} Siga el código de la Sección 6.2. Sumerja el electrodo en buffer pH 7.0, luego en pH 4.0 usando los comandos seriales \texttt{calph} para calibrar.
        \item \textbf{TDS:} Con el código de la Sección 6.1, sumerja la sonda en la solución de 1413 µS/cm y ajuste el factor de calibración según la librería.
    \end{enumerate}
    \item \textbf{Conexión y carga de código:}
    \begin{enumerate}
        \item Conecte TDS (A1), pH (A0) y MiCS-5524 (A0? ¡Cuidado! Use pines analógicos diferentes, por ejemplo TDS a A1, pH a A2, MiCS a A3). Ajuste el código.
        \item Conecte el pin EN del MiCS-5524 a D10.
        \item Cargue un código unificado que lea los tres sensores cada 2 segundos.
    \end{enumerate}
    \item \textbf{Medición de calidad de agua:}
    \begin{enumerate}
        \item Mida TDS y pH de la muestra de agua destilada (referencia).
        \item Mida agua potable de la llave.
        \item Mida agua con una cucharada de sal (alto TDS).
        \item Mida agua con vinagre (pH ácido) y agua con bicarbonato (pH básico).
        \item Registre todos los valores en una tabla.
    \end{enumerate}
    \item \textbf{Medición de calidad de aire:}
    \begin{enumerate}
        \item Espere 3 minutos a que el MiCS-5524 se caliente (LED indicador o mensaje serial).
        \item Registre los valores de CO, CH4, etanol en aire ambiente.
        \item Acerque un algodón con alcohol y observe el aumento de etanol.
        \item (Con cuidado) Acerque un encendedor sin llama (solo gas) y observe el aumento de metano/CO.
    \end{enumerate}
\end{enumerate}

\section*{Esquemático}
\begin{figure}[htbp]
\centering
\begin{circuitikz}[scale=0.7, transform shape]
\draw (0,0) node[above]{Arduino Uno};
\draw (1,0) node[ground]{GND};
\draw (1,0.5) node[anchor=south]{5V};
% pH a A2
\draw (0.5,1.5) node[anchor=east]{A2} -- (2.5,1.5);
\node at (3,1.5) {pH Sensor};
% TDS a A1
\draw (0.5,2) node[anchor=east]{A1} -- (2.5,2);
\node at (3,2) {TDS Sensor};
% MiCS-5524 a A3 y EN a D10
\draw (0.5,2.5) node[anchor=east]{A3} -- (2.5,2.5);
\draw (0.5,3) node[anchor=east]{D10} -- (2.5,3);
\node at (3,2.75) {MiCS-5524 (A0/EN)};
\end{circuitikz}
\caption{Conexión de sensores de calidad de agua (pH, TDS) y aire (MiCS-5524). Use pines analógicos distintos para evitar interferencias.}
\end{figure}

\section*{Preguntas de análisis}
\begin{enumerate}
    \item Compare los valores de TDS y pH entre agua destilada, potable y contaminada con sal. ¿Son los esperados según la teoría?
    \item ¿Por qué el agua con bicarbonato (básico) puede tener un TDS más alto que el agua neutra?
    \item ¿Qué concentración de CO detectó en aire ambiente? ¿Supera el límite de 9 ppm recomendado por la OMS?
    \item Explique por qué el MiCS-5524 requiere un período de calentamiento. ¿Qué ocurre si no se respeta?
    \item Proponga un sistema de alerta temprana usando estos tres sensores y un actuador (buzzer, relé). ¿Qué umbrales usaría?
\end{enumerate}

> Mía:
\section*{Bibliografía}
\begin{itemize}
    \item Hoja de datos MiCS-5524. SGX Sensortech.
    \item DFRobot Gravity: Analog TDS Sensor / pH Meter – Wiki oficial.
    \item WHO Guidelines for Drinking-water Quality, 4th edition.
\end{itemize}
